% asgard_boundary_conditions.tex
% Compile with: pdflatex asgard_boundary_conditions.tex  (twice, for cross-refs)
\documentclass[10pt]{article}

\usepackage[margin=0.85in]{geometry}
\usepackage{amsmath}
\usepackage{amssymb}
\usepackage{booktabs}
\usepackage{listings}
\usepackage{xcolor}
\usepackage{hyperref}
\hypersetup{colorlinks=true, linkcolor=blue!50!black, urlcolor=blue!50!black}

\newcommand{\code}[1]{\texttt{#1}}
\newcommand{\jump}[1]{[\![#1]\!]}
\newcommand{\avg}[1]{\{\!\{#1\}\!\}}
\newcommand{\dd}{\,\mathrm{d}}
\newcommand{\hF}{\widehat{F}}
\newcommand{\hq}{\widehat{q}}
\newcommand{\hf}{\widehat{f}}
\newcommand{\INF}{\textcolor{red!60!black}{\textsc{[inference]}}}
\newcommand{\UNR}{\textcolor{red!60!black}{\textsc{[unresolved]}}}
\newcommand{\fl}[2]{\texttt{#1}:#2}

\lstset{basicstyle=\ttfamily\scriptsize, language=C++,
  keywordstyle=\color{blue!45!black}, commentstyle=\color{green!30!black}\itshape,
  frame=leftline, framesep=4pt, xleftmargin=7pt, breaklines=true,
  showstringspaces=false, aboveskip=4pt, belowskip=4pt}

\setlength{\parskip}{1.5pt}
\setlength{\parindent}{0pt}
\renewcommand{\arraystretch}{1.02}

\title{\vspace{-10mm}\textbf{Boundary conditions in ASGarD}\\[0.5mm]
\normalsize How a DG weak form becomes matrix blocks, how the flags select which
blocks exist,\\ and how that composes across chains, terms and dimensions}
\author{}\date{}

\begin{document}
\maketitle
\vspace{-14mm}

{\small
\textbf{Abstract.} ASGarD has no boundary-condition module. Boundary conditions
are an emergent property of \emph{which matrix blocks get written during
assembly}, and \code{boundary\_type} is a switch selecting among three treatments
of the two unpaired edge blocks at the ends of the domain. This document derives
that architecture from the weak form, so that a reader can design boundary
conditions for \emph{new} terms rather than pattern-match existing ones. The
spine: a \code{term\_1d} chain \code{\{div, grad\}} is a genuine matrix product
$A=A_{\mathrm{div}}A_{\mathrm{grad}}$, formed at build time, with the LDG
auxiliary variable eliminated algebraically. The \code{div}'s flag sets the
numerical \emph{flux} at a wall; the \code{grad}'s flag sets the numerical
\emph{trace}. A second-order operator needs exactly one condition per wall, so at
each wall exactly one of the two links must impose its condition there.
Everything else --- the sum rule across independent terms, the emergent Robin
condition, the meaning of a Robin coefficient, and the four ways to get it wrong
--- follows. Claims are tagged \INF{} where inferred but untested and \UNR{}
where the sources did not settle them; measurements are summarised in
\S\ref{sec:verify} and recorded in full, with commands, in
\code{03\_experiments.md} beside this file.

\textbf{Source of truth.} Line numbers refer to ASGarD 0.9.1 at
\code{\textasciitilde/venvs/asgardpy/}; \code{include/} and \code{src/} are
byte-identical apart from \code{assert}$\leftrightarrow$\code{expect}. The copy
vendored in FENRIS is slightly older (\code{.cpp} lines shift by 10--20). FENRIS
solvers alias \code{bc = asgard::boundary\_type} and
\code{flux = asgard::flux\_type}; both spellings name the same enums.
}

% =========================================================================
\section{Reference card}
\label{sec:card}

\textbf{The one sign convention.} \code{term\_div(c)} assembles the DG weak form
of $+\partial_x(cf)$ into a matrix $A$, and terms are applied with
$\alpha=-1$.\footnote{\fl{asgard\_discretization.cpp}{556}, and \fl{\dots}{630}
for the implicit path.} \textbf{Throughout, $F$ is the flux in the realised
equation:}
\begin{equation}
\boxed{\;F := c\,f, \qquad \partial_t f = -\partial_x F, \qquad
       \dd N/\dd t = F(a) - F(b).\;}
\label{eq:sign}
\end{equation}
Every weak form below describes $A$, \emph{not} $\partial_t f$; the $\alpha=-1$
is paid once, here. If your PDE reads $\partial_t f=\partial_x G$ then $F=-G$ and
you pass $c=-G/f$. Cross-checked against \code{sinwav.cpp}, which solves a
right-moving wave with \code{term\_div(+1, upwind, boundary\_type::left)} plus a
\code{left\_boundary\_flux} --- correct only if \code{term\_div(+1)} contributes
$-\partial_x f$.

\begin{center}\small
\begin{tabular}{@{}p{4.4cm}p{5.1cm}p{6.0cm}@{}}
\toprule
you write & assembly effect & physical condition \\
\midrule
\code{term\_div(c, none)}      & write $\mp(c/h)\,\code{to\_*}$ into the boundary cell's diagonal & free / outflow: $\hF = c\,f^{\mathrm{int}}$ at the wall \\
\code{term\_div(c, bothsides)} & write \emph{nothing} & sealed: $\hF = 0$, exactly, at both walls \\
\code{term\_div(c, left)}      & nothing at left, free block at right & $\hF=0$ at the left wall only \\
\code{term\_div(c, periodic)}  & interior formula into the wrap blocks & wall becomes an ordinary interior edge \\
\code{term\_grad(c, none)}     & flag flips to \code{bothsides}, nothing written & trace of $f$ free \\
\code{term\_grad(c, bothsides)}& flag flips to \code{none}, block written & trace $\hf=0$ (Dirichlet) at both walls \\
\code{term\_grad(c, right)}    & block written at the \textbf{right} wall & $\hf=0$ at the \textbf{right} wall \\
\code{term\_penalty(p, bothsides)} & \emph{add} $+(p/h)\,\code{to\_*}$ & Nitsche pin $f=0$ \\
\code{set\_right\_robin(r)} (chains only) & add $+(r/h)\,\code{to\_right}$ & wall flux $\hF=rf$; needs $r\ge0$ \\
\code{+= left\_boundary\_flux(g)}  & source vector $-c\,g\,\ell_L/\sqrt{h}$ & inhomogeneous prescribed wall flux \\
\bottomrule
\end{tabular}
\end{center}

\textbf{The chain identity.} A chain is the matrix product
$A=A_{\mathrm{div}}A_{\mathrm{grad}}$, and in a dimension with mass weight $m$ it
realises
\begin{equation}
  m\,\partial_t f = -\partial_x\!\Big(\tfrac{c_d c_g}{m}\partial_x f\Big)
  \quad\Longrightarrow\quad
  \text{for physical } m\partial_t f = \partial_x(mD\,\partial_x f):\quad
  \boxed{\,c_d\,c_g = -m^2 D\,,}
\label{eq:cardcoef}
\end{equation}
so $c_d$ and $c_g$ must have \textbf{opposite signs} (\S\ref{sec:coef}).

\textbf{The flag rule.} Read each flag as the \emph{set of walls it names}:
$\code{none}\mapsto\varnothing$, $\code{left}\mapsto\{L\}$,
$\code{right}\mapsto\{R\}$, $\code{bothsides}\mapsto\{L,R\}$. The \code{div}
names walls where it seals its flux; the \code{grad} names walls where it pins
$f$. Then: \textbf{at each wall, exactly one of the two links must name that
wall.} Equivalently $\varphi_g=\mathrm{flip}(\varphi_d)$ --- the parity mnemonic,
since \code{flip} is exactly set complement.

\begin{center}\small
\begin{tabular}{@{}llll@{}}
\toprule
user flags & $\hq$ (flux) & $\hf$ (value) & realised \\
\midrule
\code{div bothsides}, \code{grad none} (default) & $0$ both & interior & \textbf{Neumann} both walls, $\dot N=0$ \\
\code{div none}, \code{grad bothsides}           & interior & $0$ both & \textbf{Dirichlet} $f=0$ both walls \\
\code{div left}, \code{grad right}               & $0$ at $L$ & $0$ at $R$ & Neumann at $L$, Dirichlet at $R$ \\
\code{div bothsides}, \code{grad bothsides}      & $0$ & $0$ & over-determined, diverges \\
\code{div none}, \code{grad none}                & interior & interior & under-determined, asymmetric \\
\bottomrule
\end{tabular}
\end{center}

\textbf{The sum rule.} Sealing a \emph{set} $S$ of terms enforces \emph{one}
condition per wall, $\sum_{k\in S}F_k\to0$ --- a \emph{signed} sum evaluated
\emph{pointwise at the wall}. Sealing drag with diffusion therefore yields an
\emph{emergent Robin} $f'/f=-A/D$: the physics of an impermeable wall, and
\emph{not} $f'=0$ (\S\ref{sec:sum}).

\textbf{Five of the eleven pitfalls in \S\ref{sec:pitfalls}.} (1)~Nothing
validates flag combinations; every pairing compiles and runs. (2)~Mass
conservation does \emph{not} catch the over-determined pair --- it conserves mass
while $f$ diverges to $1.4\times10^6$; check $\max|f|$.
(3)~\code{flux\_type::none} is silently \code{central}. (4)~Coefficients are
sampled \emph{at} the endpoints, so flags are inert wherever the Jacobian
vanishes there ($v^2$ at $v=0$, $\sin\theta$ at $\theta=0,\pi$). (5)~Same-sign
$c_d,c_g$ is an \emph{interior} instability, not a wall problem.

% =========================================================================
\section{From the weak form to matrix blocks}
\label{sec:weak}

\textbf{The broken space.} Partition $\Omega=[a,b]$ into cells
$I_j=[x_{j-1/2},x_{j+1/2}]$ of width $h$ and let $V_h=\{v:v|_{I_j}\in P^k(I_j)\}$
with \textbf{no continuity between cells}. That single fact is the origin of
everything below: a $w\in V_h$ has two values at every interior edge, $w^-$ and
$w^+$, with $\jump{w}=w^+-w^-$ and $\avg{w}=\tfrac12(w^++w^-)$. \textbf{At the
two domain edges only one trace exists; the other would come from outside the
domain. That asymmetry is the entire boundary-condition problem.}

Multiplying $\partial_x F$ by a test function $v\in P^k(I_j)$ and integrating
over one cell --- the only integration by parts in the problem --- gives a
\emph{volume} term and a \emph{surface bracket}. Since $v$ lives on $I_j$ alone,
nothing about the neighbours enters except through the endpoints:
\begin{equation}
  (A f)\big|_{I_j,\,v} = -\int_{I_j} v'\,F \dd x
        + v^-_{j+1/2}\hF_{j+1/2} - v^+_{j-1/2}\hF_{j-1/2},
  \qquad F = c\,f_h .
  \label{eq:weak}
\end{equation}
On $V_h$ the bracket is ambiguous, so DG replaces $F$ at each edge by a
\textbf{numerical flux} $\hF$: one number per edge, shared by both neighbours.
Sharing is not a convenience --- it \emph{is} the discrete statement of local
conservation. ASGarD's edge blocks give the standard upwind flux
\begin{equation}
  \hF_{j+1/2} = c\,\avg{f} - \tfrac{s|c|}{2}\jump{f}
  = \begin{cases} c f^- & c>0,\ s=+1\\ c f^+ & c<0,\ s=+1\end{cases}
  \qquad s=\code{static\_cast<int>(flux)}\in\{+1,0,-1\}.
  \label{eq:upwind}
\end{equation}
Nothing below depends on the degree $k$, the quadrature order, or on whether the
basis is nodal Legendre or hierarchical wavelet --- the wavelet transform is an
orthogonal change of basis on the \emph{same} $V_h$.

\textbf{Where $1/h$ and $2/h$ come from.} Legendre polynomials on $\xi\in[-1,1]$
are rescaled so that $\tfrac12\int P_iP_j\dd\xi=\delta_{ij}$;\footnote{\fl{asgard\_quadrature.cpp}{58--66}.}
the physical basis is $\varphi_j(x)=P_j(\xi)/\sqrt h$, making the cell mass matrix
the identity (which is why nothing inverts a mass matrix unless you asked for
one). An edge trace product is $\varphi_r(\pm1)\varphi_c(\pm1)=P_rP_c/h$ --- one
$h^{-1/2}$ per basis factor --- while $\dd\varphi_c/\dd x=(2/h)P'_c/\sqrt h$ with
$\int_{\rm cell}\dd x=(h/2)\int\dd\xi$ and the quadrature array already carrying
the $\tfrac12$. Those are exactly the two code constants,
\code{escale = 1/dx} and \code{vscale = 2/dx}.\footnote{\fl{asgard\_coefficients\_mats.hpp}{187--188}.}
No $\sqrt h$ survives: the operator is dimensionally $1/\text{length}$, as
$\partial_x$ should be.

\textbf{The four edge blocks.} Blocks are column-major with row $=$ test, column
$=$ trial. With $\ell_L=P(-1)$, $\ell_R=P(+1)$ they are rank-one outer
products,\footnote{\fl{asgard\_transformations.cpp}{55--58};
\code{gemm\_outer\_inc(n,x,y,A)} does \code{A[c*n+r] += x[r]*y[c]},
\fl{asgard\_small\_mats.hpp}{593--600}.} built once:
\begin{center}\small
\begin{tabular}{@{}llll@{}}
\toprule
\code{to\_left}$=\ell_L\!\otimes\!\ell_L$ & \code{from\_left}$=\ell_L\!\otimes\!\ell_R$ & \code{from\_right}$=\ell_R\!\otimes\!\ell_L$ & \code{to\_right}$=\ell_R\!\otimes\!\ell_R$\\
test trace at my LEFT edge $\times$ & at my LEFT edge $\times$ & at my RIGHT edge $\times$ & at my RIGHT edge $\times$ \\
\textbf{my own} value there & \textbf{left neighbour's} value & \textbf{right neighbour's} value & \textbf{my own} value there \\
\bottomrule
\end{tabular}
\end{center}
\code{to\_*} means the trace comes from \emph{this} cell (diagonal block);
\code{from\_*} from the \emph{neighbour} (off-diagonal). \textbf{These four
objects are the entire vocabulary of boundary conditions in ASGarD: every
boundary condition in the code is a decision about whether to write a multiple of
\code{to\_left} or \code{to\_right} into the first or last diagonal block.}

The operator is block-tridiagonal, with $\mathrm{lower}(0)$ and
$\mathrm{upper}(N{-}1)$ reserved as wrap blocks for
periodicity.\footnote{\fl{asgard\_block\_matrix.hpp}{317--320}; slots always
exist (\fl{asgard\_grid\_1d.hpp}{316--327}) and are zero otherwise.} For an
interior cell $i$, with $\ell=\tfrac12c(x_{i-1/2})$, $r=\tfrac12c(x_{i+1/2})$,
the assembly writes\footnote{\fl{asgard\_coefficients\_mats.hpp}{239--249}; the
volume block by \code{apply\_volume} at \fl{\dots}{210--222} as
$-c\int\varphi_r'\varphi_c$.}
$\mathrm{lower}(i)\mathrel{+}=\tfrac1h(-\ell-s|\ell|)\code{from\_left}$,
$\mathrm{diag}(i)\mathrel{+}=\tfrac1h(-\ell+s|\ell|)\code{to\_left}
+\tfrac1h(r+s|r|)\code{to\_right}$,
$\mathrm{upper}(i)\mathrel{+}=\tfrac1h(r-s|r|)\code{from\_right}$ ---
which is \eqref{eq:weak} with \eqref{eq:upwind}. For \code{upwind} and $c>0$ this
collapses to $\mathrm{lower}=-(c/h)\code{from\_left}$,
$\mathrm{diag}=+(c/h)\code{to\_right}$: the trace is taken from the \emph{left}
at every edge, which with \eqref{eq:sign} is genuine upwinding.

\subsection{Telescoping: only two blocks can break conservation}
\label{sec:telescope}

Set $v\equiv1$ in \eqref{eq:weak} --- legitimate for every $k\ge0$, since
constants lie in $P^k$ on each cell. The volume term dies because $v'=0$, and
with $\partial_t f=-Af$ the cell mass $m_j=\int_{I_j}f_h$ obeys
\begin{equation}
  \dd m_j/\dd t = \hF_{j-1/2} - \hF_{j+1/2}.
  \label{eq:ledger}
\end{equation}
Each cell is an exact ledger. Sum over $j$. Every \textbf{interior} flux appears
exactly twice --- with $-$ from cell $I_j$ and $+$ from cell $I_{j+1}$ --- and
because it is \emph{the same number} in both, the two cancel identically.
Structurally: each interior edge contributes a \emph{matched pair} of blocks, one
\code{to\_*} in cell $i$'s diagonal and one \code{from\_*} in the neighbour's
off-diagonal, and tested against $\varphi_0$ the pair cancels exactly. So
\begin{equation}
  \boxed{\ \dd/\dd t\textstyle\int_\Omega f_h \;=\; \hF_a - \hF_b\ }
  \label{eq:master}
\end{equation}
which is \eqref{eq:sign} at the discrete level. Three readings:

$\bullet$ The cancellation is \textbf{algebraic, not asymptotic}. It does not
improve with $h\to0$ or $k\to\infty$; it is exact at every resolution, for any
upwinding parameter $s$, for any coefficient $c(x)$ however rough, and for any
quadrature rule --- the volume term never entered \eqref{eq:master} at all.

$\bullet$ Therefore \textbf{the interior discretisation cannot create or destroy
mass}. It is incapable of it. Every unit of mass that appears or disappears in an
ASGarD run passed through one of exactly two numbers.

$\bullet$ Consequently, \emph{for the conservative content of a term}, a boundary
condition is a rule for those two numbers, and there is nowhere else to put one.
The qualifier is load-bearing: a second-order term has two \emph{more} wall
degrees of freedom, the traces $\hf_a,\hf_b$ of \S\ref{sec:chain}, and those do
not enter \eqref{eq:master} at all. That is exactly why the over-determined
pairing conserves mass perfectly while diverging pointwise (P2).

Only the two outermost edge blocks are unpaired. When $N>1$, cell $0$'s right
edge and cell $N{-}1$'s left edge are ordinary interior edges written
unconditionally; only the outermost two are governed by the boundary
switch.\footnote{\fl{asgard\_coefficients\_mats.hpp}{321--333};
\code{num\_cells == 1} special-cased at \fl{\dots}{263--264, 321}.}

% =========================================================================
\section{What each \code{boundary\_type} does at assembly}
\label{sec:flags}

The whole treatment is three \code{switch}es in one \code{\#pragma omp single}
block.\footnote{\fl{asgard\_coefficients\_mats.hpp}{255--351}; enum at
\fl{asgard\_pde.hpp}{109--121}. Interior cells are handled by the parallel loop
above.} At the left wall, for a \code{div}:
\begin{lstlisting}
case boundary_type::none:
case boundary_type::right:  // free on the left
  smmat::axpy(nblock, -escale * (c at xleft), basis.to_left, coeff.diag(0));  break;
default:                    // dirichlet flux, nothing to set
  break;
\end{lstlisting}
\vspace{-3mm}\footnote{\fl{asgard\_coefficients\_mats.hpp}{305--319}; the right
wall at \fl{\dots}{336--350} mirrors it, writing
\code{+escale*(c at xright)*to\_right} into \code{diag(rmost)}.}

\textbf{\code{none} --- free.} An unpaired block $\mp(c/h)\code{to\_*}$ goes into
the boundary cell's own diagonal; tested against $\varphi_0$ it is
$\hF=c\,f^{\rm int}$. The coefficient is the \emph{full} $c$, not $c/2$, and
$|c|$ never appears --- upwinding is abandoned at the wall. \textbf{This is only
well-posed on outflow.} When the characteristic points \emph{in}, extrapolating
posits a reservoir outside the domain whose density equals the solution's own
wall trace, and \eqref{eq:master} acquires $\dot N\ni-c\,f^{\rm int}(b)$: a
source \emph{proportional to the unknown}. That is an $O(1)$ modelling error, not
a discretisation error, so refinement converges faster to the wrong answer. It is
the ICRF\_2D pathology --- its drag has $c(x_{\max})<0$, so $\dot N|_b>0$:
$+63\%$ manufactured particles by $t=100$, worse under refinement. The code runs
silently; \INF{} that the absence of a guard is oversight rather than design.

\textbf{\code{bothsides} --- sealed.} The \code{default:} arm, commented
\emph{``dirichlet flux, nothing to set''}. \textbf{No \code{axpy} executes}: the
block is omitted, so $\hF=0$ identically --- enforced by \emph{absence}, hence
exact to machine precision at every level and degree.

\textbf{The ``dirichlet flux'' trap.} The comment means Dirichlet \emph{on the
flux}, $\hF=0$ --- which as a condition on $f$ is a \textbf{Neumann} condition
for a diffusive term and a no-outflow condition for an advective one. It is
\emph{not} $f=0$. Confusing the two is the commonest way to mis-flag a term.

\textbf{\code{left}/\code{right}} are \code{bothsides} on one wall only; the flag
names the wall where the flux \emph{is} pinned. \textbf{\code{periodic}} writes
the ordinary interior formula into the wrap blocks, so the domain edge becomes a
genuine interior edge joining cell $N{-}1$ to cell $0$ and conserves by the same
pairwise cancellation. \UNR{}: nothing cross-checks flags \emph{between} terms,
so mixing \code{periodic} and \code{bothsides} in one dimension across two terms
was not established.

\textbf{Sampling, and inert flags.} The wall values used above are $c$ evaluated
\textbf{exactly at the domain endpoint}.\footnote{\fl{asgard\_coefficients\_mats.hpp}{164--182}.}
So \textbf{wherever the flux coefficient vanishes at the wall, \code{none} and
\code{bothsides} assemble the same matrix} and the seal is a no-op. In FENRIS
this is routine: $c\propto v^2$ makes the $v=0$ wall inert, and
$c\propto\sin\theta$ on $[0,\pi]$ makes \emph{both} $\theta$ walls inert. Never
conclude from ``that wall is fine'' that the flags are right --- they may be
inert. Structural (from the sampling location), not separately measured.

\subsection{\code{grad}: the flag is flipped, but you get the wall you named}
\label{sec:gradflip}

Before assembly, a \code{grad}'s flag is flipped ---
\code{bothsides}$\leftrightarrow$\code{none},
\code{left}$\leftrightarrow$\code{right}, \code{periodic} untouched ``since it is
symmetric anyway''.\footnote{\fl{asgard\_coefficients\_mats.hpp}{130--148}.} A
div-style matrix $\widetilde A$ is built with the \emph{flipped} flag and then
replaced by its negative transpose with off-diagonals
swapped,\footnote{\fl{asgard\_coefficients\_mats.hpp}{355--366}, including a
final swap of the periodic wrap pair.} giving $G=-\widetilde A^{\mathsf T}$: the
discrete $\partial_x^*=-\partial_x$. Since \code{to\_left} is symmetric, a free
block $-(c/h)\code{to\_left}$ in $\widetilde A$ becomes $+(c/h)\code{to\_left}$
in $G$, so $(Gf)_r=c\int\varphi_r\partial_xf+c\varphi_r(x_0)f(x_0)+\cdots$ ---
which against the strong-form DG discretisation of $c\partial_xf$ is exactly the
wall correction with $\hf=0$. So:
\begin{quote}
For a \code{grad}, the \textbf{presence} of the boundary block pins the numerical
trace of $f$ at that wall to \textbf{zero (Dirichlet)}; its \textbf{absence}
takes the trace from the interior.
\end{quote}
Because the flag was flipped, ``block present'' corresponds to the \emph{user}
writing \code{bothsides}. \textbf{The flip exists precisely so that
\code{bothsides} means ``impose my condition here'' uniformly, whichever link you
write it on.}

\textbf{Which wall.} The internal \code{right}$\to$\code{left} flip is
bookkeeping for the div-style assembler; the negative transpose puts the block
back where you asked. Trace it: user \code{grad right} $\to$ flipped
\code{left} $\to$ left switch seals (no block), right switch writes
$+(c/h)\code{to\_right}$ $\to$ after $-\widetilde A^{\mathsf T}$ the block sits in
$\mathrm{diag}(N{-}1)$. \textbf{So \code{term\_grad(c, right)} pins $f$ at the
\emph{right} wall}, confirmed by measurement: \code{\{div bothsides, grad right\}}
diverges at $f(2)$, not $f(0)$.

A second flip, at \code{term\_1d} construction, turns a user-written
\code{upwind} on a \code{grad} into a stored
\code{downwind}.\footnote{\fl{asgard\_pde.hpp}{719--724}.} Both
\code{term\_div} and \code{term\_grad} default to \code{flux\_type::upwind} and
\code{boundary\_type::none} and take the two arguments in either
order,\footnote{\fl{asgard\_pde.hpp}{171--215}.} so writing nothing about the
flux gives \code{upwind} on both links and opposite \emph{stored} fluxes --- half
of what the interior alternation needs; see \S\ref{sec:coef} for the other half.

\subsection{\code{penalty}: \code{bothsides} adds instead of deletes}
\label{sec:penalty}

A penalty has \textbf{no volume integral}, and its interior blocks are a pure
jump form $p\jump{\varphi}\jump{f}$, so $f^{\mathsf T}Pf=+p\sum_e\jump f^2\ge0$
and with $\alpha=-1$ the contribution is
dissipative.\footnote{\fl{asgard\_coefficients\_mats.hpp}{212, 191, 120--121,
231--235}; algebraically
$\mathrm{div}_{\rm upwind}(c)=\mathrm{div}_{\rm central}(c)+(|c|/2)\cdot$penalty.}
Its boundary switch puts \code{bothsides} in the arm that \textbf{adds}
$+(p/h)\code{to\_left}$ and \code{none} in the \code{default:} that adds
nothing --- the exact opposite of the \code{div}
switch.\footnote{\fl{asgard\_coefficients\_mats.hpp}{268--279}.} The inversion
dissolves once you see \emph{what} is added: $p\varphi_r(x_0)(f(x_0)-0)$, a
one-sided jump against a \emph{zero exterior state}, i.e.\ Nitsche-style weak
enforcement of $f=0$. \textbf{For \code{div}/\code{grad} the flag names a flux to
delete; for \code{penalty} it names a value to pin.} The penalty flag is not
flipped. \code{set\_penalty(p)} on a link inherits the parent's flag, on a chain
the \emph{last} link's,\footnote{\fl{asgard\_term\_build.cpp}{838--841,
851--854, 1095--1107}.} which makes both canonical pairings self-consistent ---
but \code{div(bothsides)} plus a penalty deletes the flux block (Neumann) and
adds a Dirichlet pin, two conditions on one wall. \INF{}: over-determined like
\S\ref{sec:four}(iii); read from the dispatch, not tested.

\subsection{Inhomogeneous wall data is a source vector}
\label{sec:inhom}

\code{left\_boundary\_flux}/\code{right\_boundary\_flux} build
$-c\,g\,\ell_L/\sqrt h$ and add it to the right-hand side with the negated
weight\footnote{\fl{asgard\_term\_build.cpp}{890--931}; negation at
\fl{asgard\_term\_sources.cpp}{222, 237--239}.} --- precisely the omitted boundary
block with the unknown trace replaced by $g$. So \code{boundary\_type::left} $+$
\code{left\_boundary\_flux(g)} prescribes $\hF_a=g$, and sealing alone is the
homogeneous case ($\code{sinwav.cpp}$ is the worked example). The supplied
function must be constant in the \emph{flux} dimension but \textbf{may vary in
the others}:\footnote{\fl{asgard\_pde.hpp}{1288--1289}.} you cannot prescribe a
profile \emph{across} the wall, but you can prescribe one \emph{along} it, and
this is the only mechanism in the API for wall data depending on another
coordinate. For a chain, \code{chain\_level(d)} selects the link and the vector is
multiplied through the outer links.\footnote{\fl{asgard\_term\_build.cpp}{881--887}.
\emph{Outermost} throughout this document means the first entry in the chain
list, i.e.\ the left factor of $A_{\rm div}A_{\rm grad}$; \code{chain\_level}
defaults to $-1$, which attaches there.}

% =========================================================================
\section{Chains: the matrix product and the per-wall rule}
\label{sec:chain}

\subsection{Two wall traces, not one}

For $\partial_tf=\partial_x(D\partial_xf)$ the naive move --- integrate by parts
once and write $[vD\widehat{f'}]$ --- fails structurally: $D\partial_xf_h$ is the
derivative of a broken function, so across an edge it contains a Dirac mass the
cellwise derivative discards. One cannot integrate by parts twice on a space that
is not $H^1$. LDG refuses to, splitting into a first-order system with an
auxiliary $q$ and applying \S\ref{sec:weak}'s machinery \emph{twice}:
\begin{align}
  \textstyle\int_{I_j} w\,q_h &= -\textstyle\int_{I_j} c_g\,w'\,f_h
     + c_g\,w^-_{j+1/2}\hf_{j+1/2} - c_g\,w^+_{j-1/2}\hf_{j-1/2},
     \tag{grad}\label{eq:gradweak}\\
  (A f)\big|_{I_j,v} &= -\textstyle\int_{I_j} v'\,c_d\,q_h
     + v^-_{j+1/2}\hq_{j+1/2} - v^+_{j-1/2}\hq_{j-1/2}.
     \tag{div}\label{eq:divweak}
\end{align}
Two integrations by parts have produced \textbf{two independent numerical
traces}: $\hf$ appears \emph{only} in the \code{grad} equation, $\hq$
\emph{only} in the \code{div} equation. \eqref{eq:gradweak} has no place to put a
condition on the flux; \eqref{eq:divweak} has no place to put one on the value.
\begin{equation}
\boxed{\;
  \text{the \code{div} flag sets } \hq \Leftrightarrow \text{a condition on the
  \emph{flux}};\qquad
  \text{the \code{grad} flag sets } \hf \Leftrightarrow \text{a condition on the
  \emph{value}.}\;}
\label{eq:twotraces}
\end{equation}
They are not two chances to say the same thing; they are two \emph{different}
degrees of freedom at each wall. Note $\alpha=-1$ has \emph{not} been applied
here: \eqref{eq:gradweak}--\eqref{eq:divweak} describe $A$.

\subsection{The product, and what a chain realises}
\label{sec:coef}

\code{rebuld\_chain} [\emph{sic}] builds the \emph{last} link first and
multiplies leftwards, so the stored operator
is\footnote{\fl{asgard\_term\_build.cpp}{988--1092}; design stated at
\fl{asgard\_pde.hpp}{384--386}.} $A=A_{\rm div}A_{\rm grad}$. \textbf{$q=A_{\rm
grad}f$ is a virtual vector}: never assembled, never stored, never given its own
boundary condition. It exists only as the middle of a product collapsed at build
time. \emph{That is why the two flags interact.}

Since $A_{\rm div}\approx\partial_x(c_d\,\cdot)$, $A_{\rm grad}\approx
c_g\partial_x$, and $M^{-1}$ is applied \emph{per link} (\S\ref{sec:mass}), the
realised PDE in a dimension with mass weight $m$ is
\begin{equation}
\boxed{\;
  m\,\partial_t f = -\partial_x\!\Big(\frac{c_d c_g}{m}\,\partial_x f\Big)
  \qquad\Longrightarrow\qquad
  \text{for } m\,\partial_t f = \partial_x\big(mD\,\partial_x f\big):\quad
  c_d\,c_g = -\,m^2 D .\;}
\label{eq:coef}
\end{equation}
With no Jacobian, $c_dc_g=-D$. Two production instances: \textbf{ICRF\_1D},
$m=v^2$, $c_d=-\zeta v^2$, $c_g=+v^2$, product $-\zeta v^4=-m^2\zeta$, so
$D=\zeta$;\footnote{\code{FENRIS/ICRF\_1D/src/ICRF\_1D.cpp}:257, 276--287.} and
\textbf{LHCD\_2D} radial, $m=x^2$, $c_d=-0.25$, $c_g=x$, product
$-0.25x=-x^4\cdot\frac{1}{4x^3}$, so $D=1/(4x^3)$ --- exactly what that solver's
own comment states its collisional radial diffusivity to
be.\footnote{\code{FENRIS/LHCD\_2D/src/LHCD\_2D.cpp}:161, 175--185.}

Only the \emph{product} fixes the PDE; the split between the links is free. The
FENRIS idiom puts the Jacobian on \emph{both} links so the per-link $M^{-1}$
cancels one factor in each, leaving $M^{-1}A_{\rm
grad}\approx(1/v^2)(v^2\partial_v)=\partial_v$. What the \code{div} flag seals is
$c_dq$: the mass-weighted flux, whose wall value controls $\dd/\dd t\int mf$.

\textbf{The sign is not free, and it is the same requirement as interior
stability.} \eqref{eq:coef} with $D>0$ forces $c_d,c_g$ to have \textbf{opposite
signs}. That is not a coincidence. The LDG interior alternation --- the
\code{div} taking its trace from one side of each edge and the \code{grad} from
the other --- needs the two upwind splits $c\pm|c|$ to select opposite sides, and
the constructor's \code{upwind}$\to$\code{downwind} flip supplies only half of
that; the other half is the sign of each coefficient. Worked for constant $c$
with \code{upwind} declared on both links: a \code{div} with $c_d>0$ takes
$\hq=q^-$ and a \code{grad} with $c_g>0$ takes $\hf=f^-$ --- \emph{both from the
left}, no alternation. Making $c_d$ negative moves the div's trace to $q^+$ and
restores it. So:
\begin{quote}
\textbf{$c_d$ and $c_g$ must have opposite signs.} Getting the physical
diffusivity sign right and getting the LDG alternation right are the same
condition --- and nothing checks either.
\end{quote}
This is why P6's blow-up is possible and why it is \emph{interior}, not a wall
effect. \code{check\_chain} enforces only ``at most two non-central fluxes, never
a central mixed with a side flux''; it does not verify that the two side fluxes
point opposite ways and cannot know the sign of a variable
coefficient.\footnote{\fl{asgard\_pde.hpp}{730--745}.} \INF{}: worse,
\code{gemm\_block\_tri} computes only three bands
(\fl{asgard\_block\_matrix.hpp}{544--547}), relying on the two factors being
bidiagonal in \emph{opposite} directions; otherwise the out-of-band entries are
silently dropped and you get a consistent-looking but wrong operator. Read from
source, not measured.

Chain hygiene: identity links are stripped, a single-non-identity chain collapses
to that term, and chains of chains throw,\footnote{\fl{asgard\_pde.hpp}{524--561}.}
so \code{\{div, identity, grad\}} and \code{\{div, grad\}} build the same matrix.
Volume and identity terms have no edge blocks at all.

\subsection{The rule, stated per wall}
\label{sec:four}

Read each flag as the \textbf{set of walls it names}
($\code{none}\mapsto\varnothing$, $\code{left}\mapsto\{L\}$,
$\code{right}\mapsto\{R\}$, $\code{bothsides}\mapsto\{L,R\}$). Let $S_d$ be the
walls where the \code{div} seals its flux ($\hq=0$) and $S_g$ the walls where the
\code{grad} pins the trace ($\hf=0$). A second-order operator on an interval
needs exactly one condition per wall, so
\begin{equation}
\boxed{\;\text{at each wall, exactly one of the two links must name that wall:}
\quad S_g = \complement S_d .\;}
\label{eq:rule}
\end{equation}
Since \code{flip} is exactly set complement on this encoding, \eqref{eq:rule} is
literally $\varphi_g=\mathrm{flip}(\varphi_d)$, the parity mnemonic --- easier to
remember, but it does not tell you \emph{which} of Neumann and Dirichlet you
bought. Note \eqref{eq:rule} does \emph{not} say ``exactly one link carries
\code{bothsides}'': \code{\{div left, grad right\}} carries none and is perfectly
legal. \code{periodic} sits outside the framing
($\mathrm{flip}(\code{periodic})=\code{periodic}$) and must be matched
\code{periodic}$\leftrightarrow$\code{periodic}; see the \UNR{} in
\S\ref{sec:flags}.

The five cases are tabulated in \S\ref{sec:card}; what that table cannot say is
what each one does. \textbf{(i)~Sealed pair} ($S_d=\{L,R\}$, $S_g=\varnothing$):
the div's block is deleted, so $\dot N=0$ \emph{exactly}, while the grad still
sees the true interior trace, so $q_h$ is computed correctly right up to the
wall. Any constant is a steady state; all four FENRIS solvers use this.
\textbf{(ii)~Open pair}: mass leaves, as an absorbing wall must.
\textbf{(iii)~Both name a wall} asserts $f=0$ \emph{and} $\partial_xf=0$ at one
point; measured, divergence to $1.4\times10^6$ at level 7 growing as $h^{-3}$,
\emph{with mass conserved} (P2). \textbf{(iv)~Neither names a wall}: the outcome
is decided by the upwinding convention, which for a parabolic operator has no
physical meaning --- measured, the left wall converges to $\approx0.192$ (19\%
below the clean Neumann $0.237950$) and the right to $0$ at $O(h)$ and slightly
\emph{negative}, despite symmetric flags. \textbf{(v)~Mixed one-sided} is legal
and is the common design (\S\ref{sec:worked}); not measured, but it follows from
(i), (ii) and locality.

\textbf{Locality.} The bracket in \eqref{eq:weak} is a sum of two \emph{separate}
endpoint terms, so the walls are algebraically independent and a bad flag at one
cannot make the other diverge. Measured for \code{\{div bothsides, grad right\}}:
$f(2)$ diverges as $h^{-3}$ while $f(0)$ stays bounded and converges at $O(h)$.
\textbf{But the good wall is not ``exactly Neumann''}: it reads $0.201417$ at
level 7 (extrapolating to $\approx0.2018$) against the clean Neumann
$0.237950$ --- 15\% off, because the boundary spike redistributes the conserved
mass. Only the \emph{divergence} is local; the \emph{error} is global.

\textbf{Energy check.} Summing the two weak forms with $v=f_h$, $w=q_h$ over
cells gives the LDG energy identity, whose interior part cancels identically
under the alternating flux and whose right-wall part is
$B_b=\hf_b\hq_b-(f^--\hf_b)(q^--\hq_b)$ (full derivation in
\code{02\_weak\_form.md}\,\S7). Since nothing constrains the signs of $f^-,q^-$
independently, $B_b\le0$ for all data requires killing both products: set
\emph{one} trace to zero and the \emph{other} to the interior value --- which is
\eqref{eq:rule}, derived rather than asserted. Both consistent pairings give
$B_b=0$ and hence $\tfrac12\dd\|f_h\|^2/\dd t=-\|q_h\|^2\le0$: unconditional
dissipation at any $h$, any $k$. Both inconsistent pairings leave $\pm(fq)^-$, of
no fixed sign. The left-wall mirror carries the opposite sign, which is why case
(iv) is left--right asymmetric although its flags are symmetric --- consistent
with the measured mongrel (the argument predicts asymmetry; it does not by itself
predict the measured net mass loss).

% =========================================================================
\section{The sum rule across independent terms}
\label{sec:sum}

Real kinetic equations have several flux channels, $F=\sum_kF_k$, each assembled
as its own term with its own flag.
\begin{equation}
  \boxed{\ \text{Sealing a set } S \text{ imposes the single weak condition }
   \textstyle\sum_{k\in S}F_k \longrightarrow 0 \text{ at the wall.}\ }
  \label{eq:sumrule}
\end{equation}
\textbf{One condition per wall, not one per term.} The assembly deletes each
sealed term's block \emph{individually}, so one might expect $|S|$ conditions. It
does not.

\textbf{Derivation.} The weak form is linear in the terms, so $\hF=\sum_k\hF_k$
and \eqref{eq:master} becomes
\begin{equation}
  \dd N/\dd t = \textstyle\sum_{k\notin S}\big[\hF_k(a)-\hF_k(b)\big],
  \label{eq:leak}
\end{equation}
since sealed terms contribute nothing at the wall: \textbf{the leak rate is the
sum of the unsealed fluxes.} For the non-trivial half, apply the ledger
\eqref{eq:ledger} to the last cell $I_N$, whose right edge \emph{is} the wall:
$\dd m_N/\dd t=\sum_k\hF_k(x_{N-1/2})-\sum_{k\notin S}\hF_k(b)$. Note the
asymmetry: at the interior edge $x_{N-1/2}$ there are no flags, so \emph{all}
terms contribute; at the wall only the unsealed ones do. Now $m_N=O(h)$, so
$\dd m_N/\dd t=O(h)$, and letting $h\to0$ (the trace is continuous to
$O(h^{k+1})$ across the last cell) gives $\sum_kF_k(b)=\sum_{k\notin S}F_k(b)$,
i.e.\ \eqref{eq:sumrule}. \textbf{No steady-state assumption is needed}, but the
rule is asymptotic in $h$: exact on the numerical fluxes, $O(h^2)$ on the
post-processed pointwise $f'/f$. One equation --- the individual sealed fluxes
are \emph{not} separately zero, only their sum is. Structurally the mechanism is
trivial: all terms accumulate into one right-hand
side,\footnote{\fl{asgard\_term\_manager.cpp}{103--120}.} each open term
depositing its unpaired wall block and each sealed term depositing nothing.
\textbf{There is no per-term boundary condition anywhere in the code --- only one
flux budget per wall, and the flags choose who is in it.}

\textbf{Three regimes.} With $A_S=\sum_{k\in S}a_k(x_w)$ and
$D_S=\sum_{k\in S}D_k(x_w)$ \emph{at the wall}, \eqref{eq:sumrule} reads
$A_Sf+D_Sf'=0$, so: $D_S>0$ gives $f'/f=-A_S/D_S$ (\textbf{emergent Robin});
$D_S=0$ with $A_S\ne0$ gives $f=0$ (\textbf{Dirichlet}); $D_S=A_S=0$ gives no
condition. The first is the important one: sealing drag together with diffusion
gives $f'/f=-A/D$ --- \textbf{a Robin condition that nobody wrote}. That is the
correct physics: an impermeable wall in the presence of drag does not give
$f'=0$; it gives the drag--diffusion balance, the Maxwell--Boltzmann-like
exponential foot a no-flux kinetic equilibrium must have. Sealing gets it for
free; imposing $f'=0$ by hand would be \emph{wrong}. The second is easy to
trigger by accident: sealing a set with no $f'$ channel pins $f$, not $f'$
(measured $f(2)=1.7\times10^{-9}$ against an interior scale $4.9$). The first also
explains why sealing \emph{one} diffusion channel silences the others: $A_S=0$
forces $f'(x_w)=0$, which zeroes the wall flux of \emph{every} purely diffusive
channel there, whatever its coefficient --- measured, sealing $\{D_1,D_2\}$,
$\{D_1\}$, $\{D_2\}$ gives leak rates $9.01295/9.01295/9.01294$, identical to six
figures, although the open diffusivity is $0$, $0.5$, $1.0$.

\textbf{Signed, and strictly local.} These two properties distinguish
\eqref{eq:sumrule} from every plausible alternative, and both were measured. With
two drags of opposite sign, sealing only the negative one gives $A_S=-1$ and
predicts a wall slope of the \emph{opposite} sign, $+0.66667$: measured to
$2.8\times10^{-7}$. Tuning them to cancel gives $A_S=0$ and a pure Neumann wall:
measured $f'=3.0\times10^{-13}$. \emph{No ``first term wins'', ``largest term
wins'' or ``sum of magnitudes'' rule produces either number.} And replacing
constant $D_2=0.5$ by $D_2(x)=0.25x$ --- a different function throughout the
interior, changing the whole interior solution --- gives the \emph{same} wall
slope to 5 digits, because it has the same value \emph{at the wall}.

\textbf{Structural versus asymptotic}, and confusing them causes false alarms.
\emph{Conservation is structural}: with everything sealed, mass is conserved to
$10^{-15}$ at every level, because the block is absent from the matrix. \emph{The
pointwise wall condition is asymptotic}: measured order $2.00$--$2.02$ in $h$, so
a diagnostic checking the wall slope must expect $\approx2.5\times10^{-3}$ at
level 4 and $\approx1\times10^{-5}$ at level 8.

% =========================================================================
\section{Robin conditions}
\label{sec:robin}

The whole of \code{gen\_robin\_cmat} that touches the matrix is two
\code{axpy}s:\footnote{\fl{asgard\_coefficients\_mats.hpp}{390--394}.}
$-(r_a/h)\code{to\_left}$ into cell $0$ and $+(r_b/h)\code{to\_right}$ into the
last cell. These are \textbf{the same blocks, with the same signs and the same
$1/h$}, that the \code{none} branches of \S\ref{sec:flags} write --- with $r$ in
place of $c$. So a Robin block re-installs exactly the free/outflow wall block
with a user-chosen coefficient: $\hF^{\rm R}(a)=r_af(a)$, $\hF^{\rm
R}(b)=r_bf(b)$.

\textbf{$r$ is a flux coefficient, not a log-derivative.} Because $r$ sits where
a flux sits, $[r]=[F]/[f]$ = flux per unit $f$: for a velocity coordinate, an
\emph{effective velocity} (times whatever Jacobian the term carries). It is not
$f'/f$, whose dimension is $1/\text{length}$. Passing $A/D$ to
\code{set\_right\_robin} is dimensionally wrong, and this is not hypothetical:
ICRF\_2D's recorded history is that $r=A_{cv}/B_{cv}=2x_{\max}$ --- precisely the
log-derivative --- \emph{overshoots by ${\sim}14\times$ and the solve diverges},
while $r=A_{cv}$, the term's own wall flux coefficient, agrees with plain
\code{bothsides} to solver
tolerance.\footnote{\code{FENRIS/ICRF\_2D/src/ICRF\_2D.cpp}:610--614.}

\textbf{Equivalence to sealing.} A free \code{div} contributes $\hF_b=c(b)f^-$;
adding a Robin gives $(c(b)+r_b)f^-$. Hence
\begin{equation}
  \boxed{\ \text{free}(c) + \text{Robin}\big(r=-c(\text{wall})\big) \equiv
          \text{sealed}(c),\ }
  \label{eq:robineq}
\end{equation}
an identity of \emph{matrices}: both multiply the same rank-one block, and when
$r_b=-c(b)$ exactly the two \code{axpy}s cancel to $0$ in IEEE arithmetic.
Measured $2.1\times10^{-13}$ in a toy where $r$ was a hand-typed literal against a
quadrature-evaluated $c(b)$, and \emph{bit-identical} in ICRF\_1D and LHCD\_2D
where the two are the same expression at the same point.

\textbf{The sign condition.} The right-wall block adds
$+(r_b/h)\code{to\_right}$ to $A$, so $f^{\mathsf T}\!Af\ni r_b(f^-(b))^2$; with
$\partial_tf=-Af$,
$\tfrac12\dd\|f\|^2/\dd t \ni -r_b(f^-(b))^2 + r_a(f^+(a))^2$, giving
\begin{equation}
\boxed{\;\text{dissipative iff } r \ge 0 \text{ at the \emph{right} wall and }
r \le 0 \text{ at the \emph{left}.}\;}
\label{eq:robinsign}
\end{equation}
The mass balance agrees: $\dot N\ni r_af(a)-r_bf(b)$, so an absorbing right wall
needs $r_b>0$ and an absorbing left wall $r_a<0$. A Robin of the wrong sign is an
energy source of strength $|r|f^2$ --- a positive feedback on the wall value,
which is the free-flux-on-inflow pathology seen from the energy side. Sanity
check against \eqref{eq:robineq}: sealing a term with $c(b)<0$ is spelled
$r_b=-c(b)>0$, on the dissipative side, as it must be. Both production Robins
were of that form: ICRF\_1D's $r=+v_{\max}^2\eta(v_{\max})$ against
$c=-v^2\eta$, and LHCD\_2D's $r=+0.5$ against $c=-0.5$.

\textbf{Restrictions, all in the API.} \code{set\_left\_robin} /
\code{set\_right\_robin} take a single scalar and
\code{rassert(is\_chain())}:\footnote{\fl{asgard\_pde.hpp}{694--704}.} $r$ cannot
vary in space or time and cannot attach to a bare \code{term\_div}. It \emph{is}
mass-weighted --- the block is passed through
\code{bmass->solve}\footnote{\fl{asgard\_term\_build.cpp}{1105--1120}; likewise
the standalone \code{term\_robin} path, \fl{\dots}{861--864, 870--875}.} --- so
$r$ carries the same Jacobian weight as a \code{div}'s $c$, not the bare physical
value. And it is added \emph{after} the chain multiplication, so it acts directly
on $f$.

\textbf{Prefer sealing} for a homogeneous no-flux wall: it needs no hand
evaluation of $c$, stays exact when $c$ varies in space or time (a literal $r$ is
stale the moment the coefficient is recomputed), cannot get the sign wrong, and
composes --- realising \eqref{eq:sumrule} automatically where a hand-written
Robin would need $\sum_{k\in S}c_k(b)$ correct. A \emph{genuine} physical Robin
wall remains available: seal the \code{div} and add a Robin with the physical
(mass-weighted) $r$, so the rule reads $\sum_{k\in S}F_k(b)+rf(b)=0$ with $r$
simply one more member of the sealed set.

% =========================================================================
\section{Mass, Jacobians, and multiple dimensions}
\label{sec:mass}

\code{set\_mass} sets a positive, time-independent volume weight per dimension.
The mass matrix is a genuine Galerkin object $M_{rc}=\int\varphi_r m\varphi_c$,
Cholesky-factored once and then \textbf{inverted into every link of every
term}\footnote{\fl{asgard\_coefficients\_mats.hpp}{454--492};
\fl{asgard\_term\_build.cpp}{870--875}; \fl{asgard\_block\_matrix.cpp}{283+}. An
\code{identity} dimension gets no matrix and therefore no $M^{-1}$,
\fl{asgard\_term\_build.cpp}{683--687}.} --- so a single term stores $M^{-1}A$
and a chain stores $(M^{-1}A_0)(M^{-1}A_1)\cdots$. For one \code{term\_div(c)}
the equation actually solved is
\begin{equation}
  m(x)\,\partial_t f = -\partial_x(c f),
  \qquad\text{i.e.}\qquad \partial_t(mf) = -\partial_x(cf).
  \label{eq:masseq}
\end{equation}
\textbf{The conserved quantity is $\int mf\dd x$ and its flux is $cf$}, which is
why the coefficient handed to \code{term\_div} is effectively mass-weighted: to
represent a physical flux $\Gamma$ in a metric with weight $m$, pass
$c=m\Gamma/f$. \eqref{eq:coef} is the second-order version of the same statement.
The mass multiplies rows on the left and leaves the boundary-flux structure
untouched, so everything in \S\S\ref{sec:weak}--\ref{sec:sum} holds verbatim with
$\int f$ replaced by $\int mf$.

\textbf{A live trap in the production source.} ICRF\_1D's lambdas are \emph{named}
\code{eta\_v2} and \code{zeta\_v2} but \emph{compute} $-v^2\eta$ and
$-v^2\zeta$.\footnote{\code{FENRIS/ICRF\_1D/src/ICRF\_1D.cpp}:261--263, 277--279.}
The minus signs are essential --- they are what makes $c_dc_g=-m^2D$ come out
right --- and the names hide them. Read the body, not the name.

\textbf{Adaptivity.} Time-independent terms are assembled \emph{once at max
level} and never rebuilt,\footnote{\fl{asgard\_term\_build.cpp}{645--651, 716}.}
so the operator the solver sees is the Galerkin \emph{restriction} of a fixed
matrix. \INF{}: this is why sealing survives refinement exactly --- the level-0
constant is always in the sparse grid, so the row expressing global mass balance
is always present, and a block absent at max level is absent in every
restriction. Time-dependent terms are rebuilt at the current level, so their
$1/h$ and wall blocks change; the wall location and the seal structure do not.

\subsection{\code{term\_md}}
\label{sec:md}

A separable \code{term\_md} holds one \code{term\_1d} per dimension, each
chain-collapsed and applied as a Kronecker sweep, so the operator is $\bigotimes_d
A_d$.\footnote{\fl{asgard\_pde.hpp}{1081--1096};
\fl{asgard\_term\_build.cpp}{711--716}; \fl{asgard\_term\_manager.cpp}{67--68};
\code{term\_identity} dimensions are skipped, \fl{asgard\_kronmult.hpp}{41--46}.}

\textbf{Sealing a 1-D factor seals the whole term}, because the wall integral
factorises as (1-D wall block) $\otimes$ (mass in the other dimensions) and
zeroing the 1-D block kills the tensor product. \INF{} (structural, from the
tensor form, not separately measured): a flag on the \code{div} of dimension $d$
therefore controls dimension $d$'s wall and nothing else.

\textbf{Flux dimensions are not unique.} A \code{term\_md} carries a flux in one
dimension \emph{per \code{term\_1d} factor}; \code{flux\_dim()} returns the
\emph{first} it finds, and nothing enforces
uniqueness.\footnote{\fl{asgard\_pde.hpp}{694--712, 1264--1289}; the doxygen
claims only one is allowed, but no code checks.} The quasilinear operators in
both 2-D solvers are \code{mode::chain} \code{term\_md}s carrying fluxes in
\emph{two} dimensions at once, e.g.\
\code{term\_md(\{div\_x, diffusion, grad\_K\_xtheta\})} with the \code{div} in
$x$ and the \code{grad} in
$\theta$.\footnote{\code{FENRIS/LHCD\_2D/src/LHCD\_2D.cpp}:336--339;
\code{FENRIS/ICRF\_2D/src/ICRF\_2D.cpp}:739--742.} \INF{}: an inhomogeneous
\code{boundary\_flux} attached to such a term binds to \code{flux\_dim()}'s first
hit, silently --- a trap for chain-mode terms; read from source, not tested.

\textbf{Scope of \eqref{eq:rule}.} \code{mode::chain} links are \emph{not}
multiplied at build time; they are applied sequentially at run time through a
workspace,\footnote{\fl{asgard\_term\_manager.cpp}{103--118}.} so the intermediate
vector \emph{is} materialised and each link's wall blocks act on a real field.
\eqref{eq:rule} is derived for \code{term\_1d} chains. \UNR{}: whether it
transfers verbatim to \code{term\_md} chains was not established. What is
observable is that both production QL operators use the Neumann parity ---
\code{div bothsides} with bare \code{term\_grad} links at the default
\code{none} --- and that this is the configuration which fixed the measured
$+63\%$ particle manufacture.

\textbf{Geometric vanishing.} Combining \S\ref{sec:flags}'s endpoint sampling
with a pitch-angle Jacobian: ICRF\_2D's pitch chain is
\code{term\_div(-I(th)*sin(th), \dots, bc::bothsides)} with
\code{term\_grad(L(th)*sin(th))} on
$\theta\in[0,\pi]$.\footnote{\code{FENRIS/ICRF\_2D/src/ICRF\_2D.cpp}:510--513,
628--650.} The div coefficient vanishes at \emph{both} walls, so the block
\code{none} would write is identically zero: \code{none} and \code{bothsides}
assemble the same matrix and the seal is a no-op. Likewise $v^2$ at $v=0$. The
wall that matters is the one where the Jacobian does not vanish --- and when you
add a $\theta$-direction term whose coefficient does \emph{not} carry the
$\sin\theta$ factor, its wall block is suddenly live and the sum rule at
$\theta=0,\pi$ acquires a real member.

% =========================================================================
\section{Design recipe}
\label{sec:recipe}

\begin{enumerate}\itemsep0pt\parsep0pt
\item \textbf{Put your PDE in the convention}: $m\partial_tf=-\partial_x\sum_kF_k$.
If it reads $m\partial_tf=+\partial_x\sum_kG_k$, then $F_k=-G_k$.
\item \textbf{Read off the coefficients.} An advective channel $F=cf$ is one
\code{term\_div(c)}. A diffusive channel $F=-mD\partial_xf$ is a chain with
$c_dc_g=-m^2D$; the split is free but the signs must be \emph{opposite}. Leave
\code{flux\_type} alone --- \code{upwind} is the default on both links and the
constructor flip handles the rest.
\item \textbf{Decide the physical wall law at each wall separately}: total flux
zero (reflecting), $f$ prescribed (absorbing), $\alpha f+\beta F=0$ (partially
absorbing), or nothing (outflow). ``Total'' means the sum over \emph{all} terms;
you do not get to specify each term's flux.
\item \textbf{Choose the sealed set $S$ per wall} so $\sum_{k\in S}F_k=0$ is that
law, and have every member of $S$ name that wall on its \code{div}.
\item \textbf{Complete each chain per wall}: exactly one of its links names each
wall (\eqref{eq:rule}). If the term is in $S$ that is the \code{div}, so the
\code{grad} must \emph{not} name that wall.
\item \textbf{Audit} (\S\ref{sec:audit}).
\end{enumerate}

\subsection{Worked example, end to end}
\label{sec:worked}

$\partial_tf=\partial_v\big(A(v)f+D(v)\partial_vf\big)$ on $[0,v_{\max}]$, $D>0$,
no Jacobian ($m=1$), \textbf{perfectly reflecting at both ends}.
\emph{Step 1:} $\partial_tf=-\partial_v(F_{\rm drag}+F_{\rm diff})$ with
$F_{\rm drag}=-Af$, $F_{\rm diff}=-D\partial_vf$. \emph{Step 2:}
$c_{\rm drag}=-A(v)$; for the diffusion $c_dc_g=-D$, so take $c_d=-D(v)$,
$c_g=+1$. \emph{Steps 3--5:} reflecting means total flux zero at both walls, so
$S=\{$drag, diffusion$\}$ at both --- both \code{div}s name $\{L,R\}$ and the
\code{grad} names neither.

\begin{lstlisting}
using bc = asgard::boundary_type;                        // as in the FENRIS solvers
auto c_drag = [](std::vector<P> const &v, std::vector<P> &c) {
  for (std::size_t i = 0; i < v.size(); ++i) c[i] = -A(v[i]);      // c   = -A
};
auto c_div  = [](std::vector<P> const &v, std::vector<P> &c) {
  for (std::size_t i = 0; i < v.size(); ++i) c[i] = -D(v[i]);      // c_d = -D
};
pde += term_1d({term_div(c_drag, bc::bothsides)});                 // drag, sealed
pde += term_1d({term_div(c_div,  bc::bothsides),                   // c_d * c_g = -D
                term_grad(P{1})});                                 // c_g = +1, flag = none
\end{lstlisting}
\emph{Realised wall law}, from \eqref{eq:sumrule} at each wall: $-Af-Df'=0$,
i.e.\ $f'/f=-A/D$ --- the emergent Robin, not $f'=0$. Mass is conserved
structurally, $\dot N=0$ to round-off at every level.

\textbf{The one-line diff for absorbing at $v_{\max}$, reflecting at $v=0$.} The
walls are independent (\S\ref{sec:four}). At $L$ the flux stays sealed; at $R$ we
want $f=0$ with the flux free. So the \code{div}s name $\{L\}$ and the
\code{grad} names $\{R\}$:
\begin{lstlisting}
pde += term_1d({term_div(c_drag, bc::left)});                      // sealed at L, free at R
pde += term_1d({term_div(c_div,  bc::left),
                term_grad(P{1}, bc::right)});                      // pins f = 0 at R
\end{lstlisting}
Check \eqref{eq:rule}: $S_d=\{L\}$, $S_g=\{R\}$, complementary at both walls.
Left: reflecting, $f'/f=-A/D$, no leak. Right: $f\to0$, flux free, particles
absorbed.

\textbf{Is the free drag at the absorbing wall a problem?} Audit item 5 flags
unsealed advective terms whose characteristic points inward, and a
Fokker--Planck drag at $v_{\max}$ does point inward. It is harmless here, but only
asymptotically, and the reason is worth stating: the \code{grad}'s pin drives the
\emph{interior trace} $f^{\rm int}(v_{\max})$ to zero, and the free advective
block is $c\,f^{\rm int}$, so the destabilising flux goes to zero with it.
Measured for the pure-diffusion Dirichlet pair, the wall value converges to zero
at $O(h^3)$ ($-5.10$e--07 $\to$ $-6.37$e--08 $\to$ $-7.97$e--09 over levels
5--7). \INF{} that this suffices to tame a co-located free advective term: the
measured run has no advective channel. \textbf{Audit item 5 applies only at walls
where the value is not otherwise pinned.}

\subsection{Worked cases}

\begin{center}\small
\begin{tabular}{@{}p{2.5cm}p{2.6cm}p{4.6cm}p{5.2cm}@{}}
\toprule
physical wall law & sealed set $S$ & flags & realised condition \\
\midrule
\textbf{Reflecting} & \emph{all} flux-carrying terms &
  \code{bothsides} on every \code{div}; \code{grad}s left \code{none} &
  $\sum_kF_k=0$; with drag this is $f'/f=-A/D$, \emph{not} $f'=0$ \\
\textbf{Dirichlet} $f=0$ & $\varnothing$ &
  \code{\{div none, grad bothsides\}}; advective terms \code{none} &
  $f(\text{wall})=0$, correct Dirichlet spectrum \\
\textbf{Mixed} (reflect $L$, absorb $R$) & all terms, at $L$ only &
  \code{\{div left, grad right\}} & \S\ref{sec:worked} \\
\textbf{Prescribed inflow flux} $F=g$ & the term carrying the inflow &
  \code{bc::left} \emph{plus} \code{left\_boundary\_flux(g)} &
  $\hF_a=g$. \emph{Not} $f(a)=g$, unless the term is a bare advective \code{div} \\
\textbf{Inhomogeneous Dirichlet} $f=g$ & $\varnothing$ &
  \code{\{div none, grad bothsides\}}, boundary flux attached at the \code{grad}
  via \code{chain\_level} & \INF{}: mechanism read from source
  (\S\ref{sec:inhom}), combination untested \\
\textbf{Partially absorbing} $\alpha f+\beta F=0$ & all flux terms, plus the Robin &
  \code{bothsides} on every \code{div}, then \code{set\_right\_robin(r)}, $r\ge0$,
  mass-weighted, scalar, chain only & $\sum_{k\in S}F_k+rf=0$ \\
\textbf{Outflow} & $\varnothing$ &
  \code{none} on the advective \code{div}; keep any chain legal &
  $\hF=cf^{\rm int}$; \emph{only} valid if $c$ points out of the domain \\
\textbf{Pure Neumann} $f'=0$ (rarely wanted) & the diffusion terms only &
  \code{bothsides} on the diffusion \code{div}s, \code{none} on the drag &
  $f'=0$, and the drag then leaks --- usually a bug \\
\bottomrule
\end{tabular}
\end{center}

\subsection{The audit}
\label{sec:audit}

One row per term, one block per wall. Example at a right wall $b$ with a drag
whose $c(b)=-4$ (inward, the FENRIS-typical case):
\begin{center}\small
\begin{tabular}{@{}lllllr@{}}
\toprule
term $k$ & flux $F_k$ & $F_k(b)$ (direction) & flag & sealed? & contributes to $\dot N$ \\
\midrule
drag & $c_{\rm drag}f$ & $-4f(b)$ (inward) & \code{bothsides} & yes & --- \\
collisional diffusion & $-D_1f'$ & $-D_1f'(b)$ & \code{bothsides} & yes & --- \\
QL diffusion & $-D_2f'$ & $-D_2f'(b)$ & \code{none} & no & $+D_2f'(b)$ \\
\midrule
\multicolumn{5}{r}{\textbf{leak rate} $\dot N|_b=-\sum_{\text{unsealed}}F_k(b)$ $\rightarrow$} & $+D_2f'(b)$ \\
\bottomrule
\end{tabular}
\end{center}
\begin{enumerate}\itemsep0pt\parsep0pt
\item \textbf{The unsealed column gives the leak rate}, via \eqref{eq:leak}:
$\dot N=\sum_{k\notin S}[F_k(a)-F_k(b)]$. For an impermeable wall it must be
empty. Had the drag row been unsealed above it would contribute
$-F_{\rm drag}(b)=+4f(b)>0$: mass \emph{manufactured}, the ICRF/LHCD pathology.
\item \textbf{The sealed column gives your wall law}, and it is the only one you
get. Solve $\sum_{k\in S}F_k=0$ and check it is the physics you wanted; if
$\sum_{k\in S}D_k=0$ you have accidentally imposed Dirichlet.
\item \textbf{Per chain, per wall: exactly one link names that wall}
(\eqref{eq:rule}).
\item \textbf{Across terms, per wall: either $S$ is non-empty \emph{or} some
\code{grad} names that wall --- never both.} A \code{grad bothsides} imposes
$f=0$ there for the whole problem; combining it with a non-empty sealed set puts
two conditions on one scalar (\S\ref{sec:four}(iii), blow-up).
\item \textbf{Direction.} For every unsealed advective term at a wall where the
value is not otherwise pinned, does the characteristic leave the domain? If it
points inward, the free flux is a positive feedback on the wall value.
\item \textbf{Signs.} Per chain, $c_d$ and $c_g$ opposite (\S\ref{sec:coef}); per
Robin, $r\ge0$ at a right wall, $r\le0$ at a left one.
\item \textbf{Endpoint value.} If the coefficient vanishes there, the flag is
inert and the row is a no-op.
\end{enumerate}

\textbf{Where the recipe does not reach}, named because the failure is silent:
\textbf{(a)}~a wall law whose coefficient varies \emph{along} a wall --- seal/free
is a per-term, per-dimension binary and $r$ is a scalar, so the only mechanism
admitting dependence on another coordinate is \code{boundary\_flux}
(\S\ref{sec:inhom}); \textbf{(b)}~a genuinely non-separable wall flux, which must
be decomposed into separable \code{term\_md}s; \textbf{(c)}~cross-derivative
fluxes ($\Gamma_x$ containing $\partial_\theta f$), which are the chain-mode
\code{term\_md}s of \S\ref{sec:md}; \textbf{(d)}~translating a reflection
coefficient into $\alpha f+\beta F=0$, since the DG wall exposes only
$f(\text{wall})$ and one \emph{net} flux, never separated incoming and outgoing
halves. For (c) there is one production pattern worth copying, recorded and
measured in ICRF\_2D: the QL flux components look coupled, but because the QL
tensor is rank-1, \emph{zeroing both components} forces
$\mathbf v\!\cdot\!\nabla f=0$ and the entire QL flux drops out, leaving a purely
collisional wall law with no $\theta$ dependence. Sealing both directions can be
simpler than analysing
either.\footnote{\code{FENRIS/ICRF\_2D/src/ICRF\_2D.cpp}:585--600.}

% =========================================================================
\section{Pitfalls}
\label{sec:pitfalls}

\textbf{P1. Nothing validates flag combinations.} No check anywhere that a chain
satisfies \eqref{eq:rule}. The over-determined pair compiles, builds, runs and
blows up; the under-determined pair compiles, builds, runs and quietly produces
wrong physics.

\textbf{P2. Mass conservation does not detect the over-determined pair.} With
\code{div bothsides} the flux block is absent, so there is \emph{no mechanism} for
mass to leave, whatever the \code{grad} does. Measured at level 7: mass error
$6.3\times10^{-9}$ absolute --- $1.4\times10^{-8}$ relative to the conserved mass
$0.443$, $4\times10^{-15}$ relative to $\max|f|$ --- while $f$ diverges to
$1.4\times10^{6}$ as $h^{-3}$. The absolute error \emph{grows} $10$--$20\times$
per level ($2.9$e--11 $\to$ $5.1$e--10 $\to$ $6.3$e--09), which sharpens the point
rather than softening it. \textbf{Add a $\max|f|$ assertion.}

\textbf{P3. A one-sided bad flag corrupts the whole solution}, not just its wall:
the good wall reads $0.201417$ against the clean Neumann $0.237950$, 15\% off.

\textbf{P4. \code{flux\_type::none} is silently \code{central}.} The enumerator
has no initialiser, so it takes $\code{downwind}+1=0$ --- indistinguishable from
\code{central}. Hence $s=0$, and the dispatch calling \code{gen\_no\_flux\_cmat}
has \textbf{no \code{return}}, so execution falls through and the full
central-flux matrix is assembled on top.\footnote{\fl{asgard\_pde.hpp}{65--75};
\fl{asgard\_coefficients\_mats.hpp}{123--128, 155}.} \INF{} (strong): the no-flux
path is dead, and \code{flux\_type::none} gives a central flux \emph{with} full
edge terms, including a central-flux boundary block. Waste, not corruption (the
stale data is zeroed), and no FENRIS solver uses it.

\textbf{P5. Vanishing Jacobians make flags inert}, so a passing test at such a
wall says nothing about your flags (\S\ref{sec:flags}).

\textbf{P6. $c_d$ and $c_g$ must have opposite signs, per chain} --- equivalently
each chain's physical diffusivity $D_k=-c_dc_g/m^2$ must be $\ge0$
\emph{separately}. You cannot split a diffusion operator into a positive and a
negative chain even when the net diffusivity is positive: each chain's upwind
split $c\pm|c|$ is built from \emph{its own} coefficient and the wrong sign picks
the anti-diffusive pairing (\S\ref{sec:coef}). Measured with $D_1=+1$ fixed:
$D_2=+0.3$ behaves perfectly ($1.2\times10^{-4}$ at level 6, the expected
$O(h^2)$); $D_2=-0.05$ --- net diffusivity still $+0.95$ --- blows up by 13 orders
of magnitude in $t=0.2$; $D_2=-0.3$ reaches $\max|f|=4.3\times10^{76}$. The growth
is \textbf{interior} ($x\approx1.1$), not at a wall. State this loudly, because
the sum rule \emph{invites} the mistake: drags of opposite sign cancel across
terms perfectly well (\S\ref{sec:sum}); diffusions do not.

\textbf{P7. \code{div(bothsides)} plus a penalty is over-determined} (\INF{},
\S\ref{sec:penalty}). \quad
\textbf{P8. Free flux at an inflow wall} is a Robin of the destabilising sign and
does \emph{not} improve with refinement --- unless the wall value is pinned by
something else (\S\ref{sec:worked}). \quad
\textbf{P9. The Robin API is more restrictive than it looks}: scalar,
compile-time, chain-only, mass-weighted. \quad
\textbf{P10. \code{flux\_dim()} returns the first flux dimension, not the only
one} (\S\ref{sec:md}).

\textbf{P11. Measure wall slopes off the DG polynomial, not off a plot grid.}
Finite-differencing an interpolated grid over an interval comparable to $h$
produced apparent 21\% errors in a quantity converged to four digits (see the
corrections table). For degree 2 the solution is a quadratic on each cell, so
fitting a quadratic to 10 interior points of the last cell recovers $f$ and $f'$
at the wall to round-off.

% =========================================================================
\section{Verification}
\label{sec:verify}

Every row compares a measurement against an \emph{independently known analytic}
value. Protocol, commands, per-test tables, reproduction instructions and the
recommended regression checks are in \code{03\_experiments.md}; only the summary
is reproduced here.\footnote{One harness gotcha worth carrying over: a deck file
is read \emph{only} when passed with \code{-if}/\code{-infile}.
\code{./solver deck.txt -of out.h5} runs with every key at its default and no
error message; deck values override CLI options appearing \emph{before}
\code{-if}.} Toys are written as $\partial_tf=\partial_xG$, so in this document's
convention $F=-G$ (e.g.\ $c_{\rm drag}=-2x$); $\dd N/\dd t$ is
convention-independent and is tabulated directly.

\begin{center}\small
\begin{tabular}{@{}p{6.0cm}p{3.1cm}p{3.4cm}p{2.6cm}@{}}
\toprule
claim & predicted & measured & error \\
\midrule
\multicolumn{4}{@{}l}{\emph{Neumann pair} \code{\{div bothsides, grad default\}}} \\
eigenvalue, 4 combinations of $(L,k,D)$ & $D(k\pi/L)^2$ & agrees & $10^{-9}$--$10^{-7}$ \\
mass conservation, levels 4--8 & exact & $8.9$e--16 $\to$ $2.7$e--15 & structural \\
wall derivative under refinement & $\to0$ at $O(h^{k+1})$ & $9.7$e--08 $\to$ $2.4$e--11 & rate exactly 3 \\
\addlinespace
\multicolumn{4}{@{}l}{\emph{Dirichlet pair} \code{\{div none, grad bothsides\}}} \\
ground-mode decay rate, $L{=}2$, $D{=}1$ & $(\pi/2)^2=2.467401$ & $2.467401$ & 5.1e--09 \\
2nd, 3rd eigenvalues; $L{=}3$; $D{=}0.4$ & $\pi^2$, $(3\pi/2)^2$, $(\pi/3)^2$, $0.4(\pi/2)^2$ & all agree & $\le4.1$e--07 \\
profile relaxes to $\sin(\pi x/2)$; wall value $\to0$ & shape, then $O(h^{k+1})$ & $L_2$ dev $3.1$e--07; $5.1$e--07 $\to$ $8.0$e--09 & rate 3 \\
absolute decayed mass vs Fourier series & $0.0410189977$ & $0.0410189972$ & 1.3e--08 \\
\addlinespace
\multicolumn{4}{@{}l}{\emph{Sum rule}, $\partial_tf=\partial_x(2xf+D_1f'+D_2f')$, all 8 sealings, level 8} \\
$S{=}\{A,D_1,D_2\}$: $f'/f=-A_S/D_S$ & $-2.66667$ & $-2.66664$ & 9.3e--06 \\
$S{=}\{A,D_1\}$; $S{=}\{A,D_2\}$ & $-4.00000$; $-8.00000$ & $-3.99996$; $-7.99990$ & 1.1e--05; 1.2e--05 \\
$S{=}\{A\}$, $D_S{=}0$ $\Rightarrow$ Dirichlet & $f(2)=0$ & $1.67$e--09 & 3.4e--10 rel \\
$S$ purely diffusive $\Rightarrow$ Neumann (3 cases) & $f'=0$ & $1.4$e--07 & 1.4e--07 \\
$\dot N$ $=$ open-channel flux, all 8 sealings & --- & agrees; fully open gives $+4298$ & $\le1.7$e--04 \\
one $f'$ channel sealed silences all & identical $\dot N$ & $9.01295/9.01295/9.01294$ & 6 figures \\
wall slope under refinement, levels 4--8 & $O(h^2)$ & rates $2.00$--$2.02$; steepest case ($-8$) 3.8e--3 $\to$ 1.2e--5 & 2.5e--03 $\to$ 9.3e--06 \\
\addlinespace
\multicolumn{4}{@{}l}{\emph{Sum rule is signed and local}} \\
sealed set with $A_S=-1$ (sign flip) & $f'/f=+0.66667$ & $+0.66667$ & 2.8e--07 \\
two nonzero drags cancelling, $A_S{=}0$ & $f'=0$ (Neumann) & $3.0$e--13 & --- \\
$D_2(x){=}0.25x$ (same wall value, different interior); $D_2(x){=}0.5{+}0.75x$ ($D_S$ doubled at wall) & $-2.66667$; $-1.33333$ & $-2.66663$; $-1.33333$ & 1.3e--05; 2.0e--06 \\
\addlinespace
\multicolumn{4}{@{}l}{\emph{Robin}, $\partial_tf=\partial_x(2xf+f')$ on $[0,2]$, level 7, $t{=}4$} \\
Robin($r{=}4$) vs sealing the drag & identical matrices & max abs diff $2.1$e--13 & --- \\
sealed profile vs analytic $Ce^{-x^2}$; mass & $0.8820813908$ & $L_2$ rel $7.2$e--08; all 10 digits & --- \\
drag free with \emph{no} Robin & $0.88208$ & $1333.97$ ($1512\times$) & fails, as designed \\
\addlinespace
\multicolumn{4}{@{}l}{\emph{All six chain pairings}, $\partial_tf=\partial_xf'$, exact IC mass $0.443096843048$} \\
Neumann: mass at $t{=}1$, levels 5--7 & $0.4430968430$ & $0.4430968430$ & $\le5.0$e--16 \\
over-determined: $\max|f|$ at level 7 & --- & $1.42$e+06, grows as $h^{-3}$ & mass still exact \\
one-sided \code{grad right}: which wall & diverges at $R$ & $f(2){=}1.42$e+06, $f(0)$ bounded & --- \\
under-determined: $f(0)$; $f(2)$ & --- & $\to0.192$; $\to0^-$ at $O(h)$ & 19\% off Neumann \\
Dirichlet: mass at levels 5--7 & Fourier $0.0410189977$ & $0.0410189972$ & 1.3e--08 \\
\addlinespace
\multicolumn{4}{@{}l}{\emph{Negative-coefficient control} (P6), level 6, $t{=}0.2$, $D_1{=}+1$} \\
$D_2=+0.30$ & $f'/f=-3.076923$ & $-3.076568$ & 1.2e--04 \\
$D_2=-0.05$ (net $+0.95$); $D_2=-0.30$ & $-4.210526$; $-5.714286$ & $\max|f|=2.0$e+13; $4.3$e+76 & diverge, interior \\
\bottomrule
\end{tabular}
\end{center}

\textbf{One measured asymmetry that is \emph{not} explained.} In the Neumann
refinement study $|f'(0)|$ is consistently $2.5\times$ $|f'(2)|$, level
independently, both converging at the same $O(h^3)$ rate. An asymmetry is
expected --- the alternating flux takes $\hf$ and $\hq$ from opposite sides, so
the two walls are not equivalent --- but the \emph{value} $2.5$ is not derived
anywhere and its dependence on degree is untested. \UNR{}: treat the measured
facts (level-independent ratio, equal rates) as the claim, not the attribution.

\textbf{Corrections to earlier claims.} Three claims from an earlier session do
not reproduce as stated. None overturns the mechanism or the sum rule, but all
three would mislead someone debugging a real solver.
\begin{center}\small
\begin{tabular}{@{}p{4.4cm}p{2.5cm}p{7.6cm}@{}}
\toprule
earlier claim & status & correct statement \\
\midrule
sum-rule slopes measured $-2.60$, $-3.69$, $-6.3$ against predictions $-2.67$,
$-4$, $-8$; the last ``under-resolved'' & \textbf{measurement artifact} &
at the \emph{same} run parameters, reading $f'/f$ off the DG polynomial gives
$-2.66627$, $-3.99931$, $\mathbf{-7.99840}$ (rel.\ err $1.5$--$2.0\times10^{-4}$).
There is no boundary layer needing extra refinement. \\
\addlinespace
the over-determined pair has ``mass error $8\times10^2$'' & \textbf{wrong} &
mass is conserved ($1.4\times10^{-8}$ relative to the conserved mass). The
pathology is \emph{pointwise divergence, not mass loss}, so conservation
diagnostics do not catch it. \\
\addlinespace
with a one-sided \code{grad} flag ``the other wall stays exactly Neumann'' &
\textbf{too strong} & the \emph{divergence} is local, but the good wall reads
$0.201417$ against $0.237950$: 15\% off. The whole solution is wrong. \\
\bottomrule
\end{tabular}
\end{center}

\end{document}
